\documentclass[a4paper,11pt]{article}
\usepackage[utf8]{inputenc}
\usepackage[T1]{fontenc}
\usepackage{lmodern}
\usepackage[english]{babel}
\usepackage[colorlinks=true, urlcolor=blue]{hyperref}
\usepackage{geometry}
\geometry{left=1in, right=1in, top=0.9in, bottom=0.9in}
\usepackage{setspace}
\onehalfspacing
\usepackage{parskip}

\begin{document}

\begin{center}
{\large \textbf{Motivation Letter}} \\[4pt]
Baha Khammari \quad $\cdot$ \quad \href{mailto:b.e.khammari@gmail.com}{b.e.khammari@gmail.com} \quad $\cdot$ \quad D\"usseldorf, Germany
\end{center}

\vspace{6pt}

% --- ADAPT: Replace [Scholarship / Foundation] and tailor the
%     opening and societal-contribution paragraphs to each call. ---

\noindent Dear Members of the Selection Committee,

% === PERSONAL TRAJECTORY AND MOTIVATION ===

I apply for the [Scholarship / Programme] to support my doctoral research in applied microeconomics at [University]. Growing up bilingual in a Tunisian-German household in D\"usseldorf, I saw early that opportunity depends on the systems people are born into, not only on their effort. This observation turned into a research question during my master's degree at the University of Cologne, where I began studying whether public policies designed to equalise access to higher education actually deliver on their promise.

% === ACADEMIC ACHIEVEMENT ===

I completed my B.Sc.\ in Economics in the top 5\% of my cohort (Dean's List) and am on track to finish my M.Sc.\ by September 2026 (current GPA: 1.6/1.0). My master's thesis, graded 1.0, applied the Callaway \& Sant'Anna (2021) heterogeneity-robust staggered difference-in-differences estimator---to evaluate Brazil's ProUni scholarship programme on formal-sector wages across 558 microregions. Working with approximately five million observations of administrative employer--employee data, I found that ProUni's effects are non-monotonic and accrue disproportionately to ethnic minorities in moderate-exposure markets, while conventional estimation methods mask these patterns entirely. An Erasmus+ semester at RSM Erasmus University (Rotterdam) and two years at PwC Germany round out my profile.

% === PROPOSED RESEARCH AND ITS SIGNIFICANCE ===

My doctoral project extends this work into a systematic comparison of Brazil's three major higher education access policies---scholarships (ProUni), student loans (FIES), and affirmative-action quotas (Quota Law)---asking which design is most effective at reducing labour market inequality, and for whom. This question matters beyond Brazil: countries worldwide are expanding tertiary access through similar instruments, often without rigorous evidence on their comparative effectiveness. The project is empirically grounded in publicly available administrative data I have already worked with extensively, making timely completion realistic within a three-year PhD.

% === SOCIETAL CONTRIBUTION ===

% ADAPT this paragraph to the specific scholarship's values and mission.

My research asks whether equity policies actually produce equitable outcomes. The answer has practical implications for education policy in Germany, the EU, and the Global South. I plan to publish all data, code, and replication materials alongside the papers, so that other researchers can build on the work without having to reconstruct the dataset from scratch.

% === WHY THIS SCHOLARSHIP ===

% ADAPT: Why does this specific funding body matter to you?
% What does the scholarship enable that self-funding would not?
% (Network, mentoring, interdisciplinary exchange, freedom to
% focus full-time on research, etc.)

The [Scholarship] would let me work on this research full-time, fund conference participation and a potential research stay in Brazil, and give me access to the [Foundation's] network of [scholars / alumni / fellows]. [1--2 sentences on why this specific foundation matters to you.]

\vspace{8pt}
\noindent Yours sincerely, \\
Baha Khammari

\end{document}
